
% ============================================================================
\clearpage
\section*{v4.10.0-r17 Source-Only: Status Notes on the Clock Functional, and on the Counting Rule as an Irreducibility Condition}
\addcontentsline{toc}{section}{v4.10.0-r17 Source-Only: Status Notes on the Clock Functional, and on the Counting Rule as an Irreducibility Condition}
\label{sec:v4100-r17-clock-irreducibility}
% ============================================================================

\paragraph{Purpose.}
This layer records the explorations of 15 November 2026 that followed v4.10.0-r16: two OP3 items (the photon-clock functional
and a reading of the overlap \(\tfrac12\)) and four OP1 items (the counting rule per frame as an irreducibility condition, the
Barbero connection as a point of a doublet, Penrose strand bundles, and the withdrawal of a motivation given during the
work).  All items are explorations without pre-registration and without comparison with data; the counting pair of H98, its
values, the Koide value and the lepton phase were known in advance.  The notes are
\texttt{notes/\allowbreak research/\allowbreak v4\_10\_0\_r17/} and the folders cited there.  Earlier layers are not
rewritten.

\noindent\textbf{Source-only status.}
This v4.10.0-r17 layer is source-only.  It does not build a PDF, does not assign a new DOI,
and does not authorize a public release.  The v5.1.0 Version DOI
\href{https://doi.org/10.5281/zenodo.24237112}{10.5281/zenodo.24237112} and the Concept DOI
\href{https://doi.org/10.5281/zenodo.19210728}{10.5281/zenodo.19210728} are unchanged by it.

\subsection*{v4.10.0-r17 status notes}
\label{subsec:v4100-r17-status}

\paragraph{Theoremlet v4.10.0-r17.1 (the Koide condition has the form of the photon-clock functional).}
The photon-clock probe of v4.7.0-r90 is \(1+\eta\,|\sum_mw_me^{i\varphi_m}|^2/(\sum_m|w_m|)^2\); a window \(V\) keeps its selector while the
coherence functional is at least \(\rho_V\) (\(\rho_{V3}=\tfrac35\), \(\rho_{V5}=\tfrac25\), \(\rho_{V8}=\tfrac79\)).  The Koide condition of r16.9(b) is exactly
this functional with weights \(\sqrt{m_k}/\sum\sqrt m\) and phases \(-2\pi/3,0,2\pi/3\): value \(0.499991\).  No window threshold lies at
\(\tfrac12\); the drift of \(120^\circ\) per readout lies far outside the tolerance \(26.57^\circ\) of v4.9.0-r12.2; and identifying clock
frequencies with the field frequency would require the identification the source declines (v4.9.0-r12) and would touch the
forbidden global photon-clock statement.  The clock does not fix \(\sum p^2=\tfrac23\).

\paragraph{Remark v4.10.0-r17.2 (the overlap \(\tfrac12\) as maximal uncertainty; reading).}
Equivalent readings of \(K=\tfrac23\) are on record since 13 November 2026 (purity \(\tfrac23\) of the normalised Gram matrix, middle of the
range \([\tfrac13,1]\), equal weight of diagonal and off-diagonal, pairwise fidelity \(\tfrac12\)).  Whether one readout confirms another is
one bit with probability \(|\langle k|l\rangle|^2\); its entropy is maximal at \(\tfrac12\).  The rule ``different readouts of one object leave
maximally open whether they show the same'' gives the octant of r16.7 and \(K=\tfrac23\) in the time cycle of r16.9 from one source.  The
maximisation is assumed; no new test follows (neutrinos are covered by H2, H99, H100; quarks would need an exception).

\paragraph{Theoremlet v4.10.0-r17.3 (the frame maps act irreducibly only on \(j=\tfrac12\)).}
Assume (A) that the frame maps act on a horizon puncture of spin \(j\) through the same \(SU(2)\) as on the doublet \(\chi\), that is as
\(\exp(-i\pi J_k)\) about three orthogonal axes (r16.7).  Then, checked for \(j=\tfrac12\) to \(20\):
\begin{enumerate}[leftmargin=2em]
\item[(a)] for half-integer \(j\) the maps anticommute and the space splits into \((2j+1)/2\) copies of the two-dimensional
representation; for integer \(j\) they commute and the space splits into characters; they act irreducibly only for \(j=\tfrac12\);
\item[(b)] the extremal doublet \(|j,\pm j\rangle\) is invariant for every \(j\) (as are all pairs \(|j,\pm m\rangle\));
\item[(c)] the rules ``T counts punctures on which the frame maps act irreducibly'' (\(j=\tfrac12\), weight \(2\)) and ``M counts the
invariant extremal doublet per \(j\)'' (all \(j\), weight \(2\)) give \(\gamma_M=0.237533\), \(\gamma_T=0.127384\), \(\gamma_L=1.864700\),
\(\beta_{\rm sym}=0.54941\), the pair of H98.  Controls: Kramers for T gives \(\gamma_L=1.2559\); the weight \(2j+1\) for M gives \(2.1515\).
\end{enumerate}
Reading: the energy condition P2 of r09.2, withdrawn in r10.2, is replaced by a group-theoretic readout condition, a restricted
readout in the sense of r10.4(c).  M still needs P1 (extremal).  (A) and the asymmetry between M and T are assumed.

\paragraph{Theoremlet v4.10.0-r17.4 (the Barbero connection as a point of the doublet \((A_+,A_-)\)).}
With \(A_\pm=\Gamma\pm iK\), \(A_\gamma=\Gamma+\gamma K=c_+A_++c_-A_-\), \(c_\pm=(1\mp i\gamma)/2\).  Real \(\gamma\) lies on the equator of the doublet
(\(|c_+|=|c_-|\)), \(\gamma=\pm i\) at the poles, the same geometry as \(\chi\).  Acting as on \(\chi\), the frame maps send \(\gamma\) to \(-1/\gamma\), \(-\gamma\),
\(1/\gamma\); none maps \(\gamma_M\) to \(\gamma_T\).  The six poles of r16.7 correspond to \(\gamma=0\) or \(\infty\), \(\pm1\), \(\pm i\).  This acts on the
parameter \(\gamma\), not on puncture states, so it is not the bridge (A); the frame difference sits in the counting.  Seen afterwards
and not used: \(2\arctan\gamma_M=26.72^\circ\) lies near \(\arctan\tfrac12=26.57^\circ\).

\paragraph{Theoremlet v4.10.0-r17.5 (Penrose strand bundles and statistics).}
\begin{enumerate}[leftmargin=2em]
\item[(a)] On \(\mathrm{Sym}^{2j}(\mathbb C^2)\), \(\sigma_k^{\otimes2j}\) equals \(\exp(-i\pi J_k)\) up to a phase (checked for \(2j=1\) to \(8\)).  If a spin-\(j\) edge is a
symmetrised bundle of \(2j\) two-component strands and a strand is a quantum of \(\chi\), the frame maps act on the puncture as (A)
assumes.
\item[(b)] Antisymmetric powers of a doublet: \(\Lambda^1\) is the doublet (\(j=\tfrac12\)), \(\Lambda^2\) a singlet (\(j=0\), no area), \(\Lambda^n=0\) for \(n\ge3\).
Bosonic strands give all \(j\) (the M rule), fermionic strands only \(j=\tfrac12\) (the T rule).
\item[(c)] The aligned bundles \(|{\uparrow}\rangle^{\otimes2j}\), \(|{\downarrow}\rangle^{\otimes2j}\) are \(|j,\pm j\rangle\), the extremal doublet of M.
\end{enumerate}
Compared with the guards of CT-004 (register unchanged): ``T = string sector'' would no longer be needed and ``Chan--Paton'' would
be replaced by ``a strand is a quantum of \(\chi\)''; the guard on statistics loses its tension of content, since T11 uses bosonic and
fermionic for the \(J\)-parity of one object, not for the exchange of strands; the guard on the spectrum (H104) is unchanged.  New
assumptions: (S1) a strand is a quantum of \(\chi\); (S2) T reads strands fermionically, M bosonically; (S3) M counts only aligned
bundles.

\paragraph{Correction v4.10.0-r17.6 (a motivation for (S2) withdrawn).}
During the work (S2) was motivated by \(R^2=-1\) on the doublet (r14.3(d)), the sign of a fermion under a full turn.  This is
withdrawn.  \(R^2\) (\(z\mapsto-z=C=PT\), v4.9.0-r02) acts on a bundle of \(n\) strands as \((-1)^n\); the condition \(R^2=-1\) keeps all
half-integer \(j\), the Kramers rule, which is excluded.  The three frame maps have determinant \(-1\) on \(\chi\), act as \(-1\) on the
singlet and commute with the exchange of two strands; they do not see the statistics.  (S2) is an assumption without a
reason in the frame structure.

\subsection*{v4.10.0-r17 status in the three frames}
\label{subsec:v4100-r17-frames}

\begin{itemize}[leftmargin=2em]
\item M: counts at every \(j\) the invariant extremal doublet (r17.3(c)), the aligned bundles of bosonic strands (r17.5).
\item T: counts only punctures on which the frame maps act irreducibly, \(j=\tfrac12\) (r17.3), equivalently single fermionic strands
(r17.5(b)); the reason for the fermionic reading is open (r17.6).
\item P: not evaluated.
\end{itemize}

\subsection*{v4.10.0-r17 register entries}
\label{subsec:v4100-r17-register}

\begin{center}
\small
\begin{tabular}{@{}>{\raggedright\arraybackslash}p{0.34\textwidth}>{\raggedright\arraybackslash}p{0.58\textwidth}@{}}
\toprule
statement & status after v4.10.0-r17\\
\midrule
photon clock and Koide & same functional; no threshold at \(\tfrac12\); no identification of frequencies (r17.1)\\
overlap \(\tfrac12\) as maximal uncertainty & reading; maximisation assumed (r17.2)\\
counting rule per frame & H98 pair from ``T irreducible, M invariant extremal doublet'' under (A) (r17.3)\\
Barbero connection & real \(\gamma\) on the equator of \((A_+,A_-)\); frame maps act on \(\gamma\), not the bridge (r17.4)\\
strand bundles & give (A) if a strand is a quantum of \(\chi\); Bose gives M, Fermi gives T (r17.5)\\
\(R^2=-1\) as reason for (S2) & withdrawn; it gives Kramers (r17.6)\\
CT-001 to CT-004 & unchanged\\
H100 / H102 & pending / pending\\
\bottomrule
\end{tabular}
\end{center}

\subsection*{v4.10.0-r17 ledger and guard update}
\label{subsec:v4100-r17-ledger-guard}

\[
R_{\rm chk}^{v4100\text{-}r17}=(1,1,0)
\]
for (status notes recorded; formulas checked in the cited scripts; (A), (S1)--(S3), the overlap \(\tfrac12\) and the core size not
derived).

The layer spends no guard:
\[
R_{\rm glob}^{(v4100\text{-}r17,+)}=(0,1,1,1,1),
\qquad
\Gamma_{\rm guard}^{v4100\text{-}r17}=(1,1,1,1,1),
\qquad
\beta_4=1/256\ \text{unspent}.
\]
OP~1, OP~3 and OP~4 remain open.  The layer does not imply exact finite \(C5\) zero closure, exact all-readout terminality,
\(C6/T2\), a global photon-clock/RH theorem, or global all-readout closure.
